\documentclass{article}
\author{Dawid S. Jakubczyk, Jakub Daszewski}
\usepackage{polski}
\usepackage[table]{xcolor}
\usepackage[a4paper, margin=1in]{geometry}
\title{Model biznesowy: Przychodnia medyczna}
\begin{document}
	\begin{tabular}{|m{10em}|m{20em}|}
		\hline
		Nazwa &Wizyta lekarska \\
		\hline
		Numer PU & 2\\
		\hline
		Aktorzy & Pacjent, Lekarz\\
		\hline
		Krotki opis & pacjent konsultuje sie z lekarzem\\
		\hline
		Warunki wstepne & wizyta jest ustalona w terminarzu\\
		\hline
		Warunki koncowe & lekarz wpisal dane o wizycie do systemu\\
	
		\hline
		A
		Glowny przepływ zdarzeń& 1.Pacjent przychodzi na wizyte \newline
								 2.nastepuje konsultacja lekarska \newline
								 3. Wizyta sie konczy  \\
		\hline
		Alternatywny przeplyw zdarzen &  2b. lekarz wystawia skierowanie\\
		\hline
		Specjalne wymagania & brak\\
		\hline
		Notatki& brak\\
		\hline
	\end{tabular}
	\paragraph{}
	\begin{tabular}{|m{10em}|m{20em}|}
		\hline		
		Nazwa &Rejestracja Pacjenta \\
		\hline
		Numer PU & 1\\
		\hline
		Aktorzy & Pacjent, pracownik rejestracji\\
		\hline
		Krotki opis & Pracownik wpisuje pacjenta i ustala dane rezerwacji\\
		\hline
		Warunki wstepne & Co najmniej jeden lekarz i pokoj musi byc dostepny\\
		\hline
		Warunki koncowe & System odnotowuje pokoj lekarza jako zajete i aktualizuje harmonogram\\
		\hline
		Glowny przepływ zdarzeń& 1.Pracownik dostaje informacje od Pacjenta o terminie dostepnosci pacjenta 
		\newline 2. Pracownik wybiera dodaj wizyte i wpisuje dostepnosc pacjenta 
		\newline 3.System wybiera i rezerwuje termin 
		\newline 4.System potwerdza rezerwacje komunikatem "zarejestrowano"\\
		\hline
		Alternatywny przeplyw zdarzen &  2a System wyswietla komunikat brak pasujacego do kryteriow terminu\\
		\hline
		Specjalne wymagania & Czas rejestracji powinien byc krotszy niz 3 sekundy\\
		\hline
		Notatki& brak\\
		\hline
	\end{tabular}
\end{document}